\section{模型假设}

题目对若干表述留有解释余地，而这些口径直接决定结果的含义。本文在建模阶段即锁定基准口径，并在每条假设后说明备选解释的影响。

\begin{enumerate}
  \item 每个 10 分钟时段内的负荷功率、光伏功率与电价取该时段的平均值，时段电量为功率乘以时段长度 $\Delta t = 1/6$ 小时；
  \item 附件 1 的时刻标签按所在时段的终点解释，并按其行先后顺序与结果模板的 144 个时段一一对应。若把同一时刻改按区间起点解释，全部结果将整体平移一个时段，但各时段的相对调度结构与全天购电量、购电费不变；
  \item 储能设备的充电效率与放电效率分别为 90\%，即充电时只有 90\% 的电能进入储能，放电时需消耗 $1/0.9$ 倍的储电量才能送出 1 单位电能。若把 90\% 理解为一个充放循环的往返效率，则单程效率约为 $\sqrt{0.9}$，储能可套利的电量会相应减少；
  \item 同一时段不同时充电和放电；
  \item 储能设备不允许向外部电网售电，光伏出力在满足负荷与充电需求后仍有富余时只能弃光，弃光不产生收益，也不计入购电费用；
  \item 问题三的调整购电采用净结算口径：计划购电量高于调整购电量的部分按 50\% 电价承担违约费用，调整购电量高于计划购电量的部分按 1.5 倍电价计费。若改为``计划量全额付费并另加 50\% 违约费''，下调的边际成本将由 0.5 倍电价升至 1.5 倍电价，滚动调整会明显减少，日内更新的价值被低估；
  \item 问题四的实时电价采用严格因果口径：现实策略只能使用决策时刻以前已经知道的电价，附件 4 中当天的未来实际电价仅用于事后结算与价格完全信息基准。若假定当天全天电价已知，则现实策略即等同于完全信息基准，全年费用会相应下移。
\end{enumerate}